% !TEX TS-program = pdflatex
% !TeX program = pdflatex
% !TEX encoding = UTF-8
% !TEX spellcheck = en

\documentclass[11pt, a4paper]{article}
%\usepackage{fullpage}
\usepackage[left=1.5cm,right=1.5cm,top=1cm,bottom=2cm]{geometry}
\usepackage[fleqn]{amsmath}
\usepackage{amssymb}
%\usepackage{indentfirst}
\usepackage[T1]{fontenc}
\usepackage[utf8]{inputenc}
\usepackage[french,english]{babel}
\usepackage{txfonts}
\usepackage[]{graphicx}
\usepackage{multirow}
\usepackage{hyperref}
\usepackage{parskip}
\usepackage{multicol}
\usepackage{wrapfig}
\usepackage{fancyhdr}
\usepackage{hyperref}
\usepackage{tcolorbox}
\usepackage{natbib}

\usepackage{turnstile}%Induction symbole

\renewcommand{\baselinestretch}{1}

\setlength{\parindent}{24pt}

%\usepackage{etoolbox}
%\patchcmd{\thebibliography}{\section*{\refname}}{}{}{}

\fancyhf{}

\lfoot{Abdelkrime Aries}
\cfoot{2023-2024/ESI/2CSSID/NLP}
\rfoot{\textbf{\thepage}}
%
\renewcommand{\headrulewidth}{0pt} % remove lines as well
\renewcommand{\footrulewidth}{1pt}

\newcommand\repeatstr[1]{\leavevmode\xleaders\hbox{#1}\hfill\kern0pt}

\renewcommand{\bibsection}{}
\bibliographystyle{humannat}

\newcommand{\kurl}[1]{{\scriptsize\bfseries\url{#1}}}

\setlength{\parindent}{0pt}

\begin{document}

%\selectlanguage {french}
%\pagestyle{empty} 
\pagestyle{fancy}

\noindent
\begin{tabular}{ll}
\multirow{3}{*}{\includegraphics[width=1.6cm]{../logo/esi.nlp.pdf}} & \'Ecole national Supérieure d'Informatique, Algiers\\
& 2CSSID (2023-2024)\\
& Natural Language Processing (NLP)
\end{tabular}\\[.25cm]
\noindent\rule{\textwidth}{1pt}\\[-0.25cm]
\begin{center}
{\LARGE \textbf{Syllabus}}
\begin{flushright}
	Dr. Abdelkrime Aries
\end{flushright}
\end{center}\vspace*{-0.25cm}
\noindent\rule{\textwidth}{1pt}

\begin{abstract}
	Natural Language Processing (NLP) is an interdisciplinary field involving linguistics, computer science, and artificial intelligence. 
	It is a domain of research and computer application dedicated to the understanding and generation of human language by computers. 
	With the emergence of language models such as OpenAI's GPT, Google's BERT, and others, NLP has become a rapidly growing field with increasing importance in various industries and domains.
	
	NLP covers a wide range of topics, and it would be impossible to cover them all in a week. 
	However, in this seminar, we will focus on the most relevant and current tasks in scientific research.
\end{abstract}

\section{General information}

\begin{minipage}{0.44\textwidth}
		
	\subsection{Prerequisites}
	
	\begin{itemize}
		\item Basic knowledge of languages (vocabulary, grammar, etc.);
		\item Data structures (1CP-ALSDD);
		\item Probability and statistics \\(2CP-PRST1, 2CP-PRST2);
		\item Algebra and calculus \\(2CP-ANAL3, 2CP-ALG3, 2CP-ANAL4);
		\item Language theory (1CS-THP);
		\item Operations research (1CS-RO);
		\item Machine learning (2CS-ML or 2CS-BDM).
	\end{itemize}

\end{minipage}
\begin{minipage}{0.05\textwidth}
	\hspace{2pt}
\end{minipage}
\begin{minipage}{0.49\textwidth}
	\subsection{Objectives}
	
	\begin{itemize}
		\item Learn some linguistic concepts;
		\item Apply concepts covered in the prerequisite modules to solve NLP tasks;
		\item Discover some NLP tools and resources;
		\item Understand the research directions in this field;
		\item Implement simple solutions to a few problems;
		\item Know how to apply the concepts of language theory to process natural languages;
		\item Get an idea of some of the work carried out within ESI in this area.
	\end{itemize}
	
\end{minipage}

\subsection{Office hours}

If there are questions about a lecture, lab, workshop, tutorial or anything related to NLP, you can join me following these recommendations:
\begin{itemize}
	\item \textbf{Location}: New laboratory building (near DG), LCSI lab (first floor), office N 68 (left).
	\item \textbf{Time}: Each Tuesday from 13h30 to 16h30.
\end{itemize}

\subsection{Tools}

Labs do not need machine power, since our objective is to understand concepts and not to start a big project. 
In this case, students can bring their own machines (at least one by a team). 
Workshops need some machine power, so you can use your machines if they are powerful; otherwise, you can use online platforms such as Google colab.
More information will be presented in both "lab" and "workshop" sections.

\section{Outline}

In this section, we will present different course components.

\subsection{Lectures}

Their goal is to understand different concepts and solutions to NLP problems. 
Each lecture starts by presenting the problem in hand. 
Then, a short description of linguistic background is shown.
Finally, different approaches are presented to automatically solve the problem.
These are all the chapters:
\begin{enumerate}
	\item Introduction to NLP;
	\item ML for NLP;
	\item Basic text processing;
	\item Language models;
	\item Part of Speech (PoS) tagging;
	\item Parsing;
	\item Word semantics;
	\item Sentence meaning;
	\item Coreference detection;
	\item Discourse coherence;
	\item Some applications.
\end{enumerate}

The lecture is designed to be interactive; students can ask questions and discuss what is being presented during the presentation. 
Of course, the discussion must be organized (no need to jump into another subject; you have to discuss the subject in hand).
If you did not understand, you can always ask for clarification for a very specific aspect ("\textit{I did not understand all}" is very broad). 
Avoid asking too many simple questions to which you already have the answer (\textcolor{red}{I am talking to you DAOUD and all your clones :)})

\subsection{Demonstrations}

There are some demonstrations of some well known tools in the realm of NLP. 
Unfortunately, they were not updated since 2022, so some demonstrations may be outdated (especially for APIs with a low life cycle).
They are still on the old repository (the French one):
\begin{itemize}
	\item \url{https://github.com/projeduc/ESI_TALN/tree/master/tuto}
\end{itemize}

\subsection{Tutorials}

These are some exercises to solve, which are assisted by the professor.
You know them as "TD". 
Not graded, although they are essential for understanding.
This is an example:
\begin{itemize}
	\item \url{https://github.com/projeduc/ESI_NLP/releases/tag/0.15}
\end{itemize}

\subsection{Quizzes}

At the end of each chapter, you will receive a quiz.
The goal is to encourage course review. 
This will prepare students for the exam (at least the knowledge part).
It is an individual work.
So, those who work together or pass the answers to their colleagues are dooming them to have lower grades in the written exams.
This is because the goal of this course component was not reached.
%
Anyway, each quiz will not take more than 20 minutes.
For those who attend and listen, it will be less.
You have \textbf{almost} one week to complete it. 
\textcolor{red}{I said almost, which means less than a week (I did not receive the schedule yet)}.
The quizzes are \textbf{graded}.

\subsection{Labs}

Labs are used to promote understanding the theoretical notions and the theory/practice transition.
This can be done by programming a concept from scratch (We use Python).
Since 1 hour is not enough to code an entire system, we will use Gaps filling.
The system is coded by the professor, then some parts are deleted.
Students must complete the system and answer some additional questions about the system.

There will be 4 labs which are \textbf{graded}.
They are face-to-face; so, you have to attend.
The labs' schedule will be shared afterwards.
Here some previous examples:
\begin{itemize}
	\item TP01 fouille des contacts: \url{https://github.com/projeduc/ESI_TALN/releases/tag/3.1}
	\item TP02 autocompletion: \url{https://github.com/projeduc/ESI_TALN/releases/tag/3.2}
	\item TP03 analyse morpho-syntaxique: \url{https://github.com/projeduc/ESI_TALN/releases/tag/3.4}
	\item TP04 Analyse syntaxique avec CKY: \url{https://github.com/projeduc/ESI_TALN/releases/tag/3.5}
	\item Lab01 Naive Spell Checker: \url{https://github.com/projeduc/ESI_NLP/releases/tag/0.12}
	\item Lab02 Natural Language detector: \url{https://github.com/projeduc/ESI_NLP/releases/tag/0.13}
	\item Lab03 Parsing: \url{https://github.com/projeduc/ESI_NLP/releases/tag/0.14}
	\item Lab04 Plagiarism detection: \url{https://github.com/projeduc/ESI_NLP/releases/tag/0.16}
\end{itemize}

\subsection{Workshops}

Workshops are practicals which seek to understand how to use high-level APIs to construct an application.
They are not graded.
The tools are those presented in the demonstration subsection.
Here are some examples:
\begin{itemize}
	\item Jugements d'acceptabilité: \url{https://github.com/projeduc/ESI_TALN/releases/tag/3.3}
	\item Sentiment analysis: \url{https://github.com/projeduc/ESI_NLP/releases/tag/0.17}
\end{itemize}

\section{Evaluation policy}

In this section, we will discover how students will be evaluated.

\subsection{Grading}

Students will be graded based on 4 evaluations:
\begin{itemize}
	\item \textbf{Quizzes}: 10\% of the quizzes' average grade (see "quizzes" subsection). 
	A quiz evaluates how well a student understands and reviews a chapter. 
	It also evaluates 
	\item \textbf{Labs}: 20\% of the labs' average grade (see "labs" subsection).
	A lab evaluates how 
	\item \textbf{Midterm Exam}: 30\%. It is a written exam (no PC, no documents allowed). 
	Usually, all chapters from the first one till "parsing".
	\item \textbf{Final Exam}: 40\%. The same thing as the midterm. The chapters after parsing, although some concepts from first chapters may be used such as those of pre-processing.
\end{itemize}

Midterm and final exams, both seek to assess a student's capacity to understand and use different concepts in a limited time. 
So, it is not enough to master a subject; you must know how to manage your time as well.
To be fair, these exams evaluate three skills:
\begin{itemize}
	\item \textbf{Knowledge test}: mostly a quiz to test a student's knowledge. 
	Knowing things is an important aspect.
	Even if we have access to information, we must have the ability to filter right information based on what we already know.
	\item \textbf{Application}: the ability to apply algorithms or execute calculations.
	Even if this aspect can be automated, we must know how it is done.
	In the end, an engineer must understand how a concept functions so it can be implemented or enhanced.
	\item \textbf{Critical Thinking Test}: the ability to propose a solution to some problem we did not see using concepts we already know.
	This test can be a problem to solve or a solution to criticize.
\end{itemize}

\subsection{Plagiarism}

Usually, plagiarism happens in graded labs. 
Giving your solution to another person means both will get 0 in similar parts. 
Although there is no restriction on using external resources, it is not advised to use a code from "stackoverflow", past labs, chatGPT generated, etc. 
If you use external resources, you have to adapt their content to the lab which is mostly specific and guided so no external code will be embedded easily. 
In this spirit, when two teams have the same code, it will be considered as plagiarism even if their is no direct sharing between them. 

\subsection{Missed exams}

This is the policy when a student fail to attend one or more evaluation.
\begin{itemize}
	\item \textbf{Labs}: The labs must be face-to-face one hour evaluation. 
	\begin{itemize}
		\item If you had an injury/illness, you have to inform me before the session. You will be granted an online session. In this case, you will work alone and so are your teammate. You are required to afford a certificate; otherwise it is a zero.
		\item If you have a more serious case such as the death of a family number or an illness which prevent you from even working from home, you have to inform me beforehand. A justification must be afforded afterwards. If the case is justified, the replacement method will be decided afterwards based on the case.
		\item If you miss either Midterm or final exam, you must attend a replacement exam. The latter will involve \textbf{ALL} chapters since it can replace both. 
		Even if there is one person who did not attend one of the two exams, the replacement still covers all chapters.
		\item You cannot miss both Midterm or final exam. If so, the replacement will be considered as a final exam (being the one with higher percentage) and the midterm will be 0.
	\end{itemize}
\end{itemize}


\section{Bibliography}

You have to download the textbooks in subsection \ref{sec:txtbk}.

\subsection{Similar courses}

\begin{itemize}
	\item "\textit{NLP S21}" (up to date).
	Carnegie Mellon Univ. 
	\kurl{http://demo.clab.cs.cmu.edu/NLP/}
	
	\item "\textit{CS388: NLP}" (2018). 
	Univ. of Texas. 
	\kurl{https://www.cs.utexas.edu/~mooney/cs388/}
	
	\item "\textit{CSE 517: NLP}" (NA).
	Univ. of Washington.
	\kurl{https://courses.cs.washington.edu/courses/cse517/}
	
	\item "\textit{CS 294-5: Statistical NLP}" (2005). 
	Univ. of Berkely. 
	\kurl{https://people.eecs.berkeley.edu/~klein/cs294-5/index.html}
	
	\item "\textit{CS224n: NLP with Deep Learning}" (up to date).
	Stanford Univ.
	\kurl{https://web.stanford.edu/class/cs224n/}
	
	\item "\textit{CS 5340/6340: NLP}" (up to date). 
	Univ. of Utah.
	\kurl{https://my.eng.utah.edu/~cs5340/schedule.html}
\end{itemize}

\subsection{To read (NLP)}
\label{sec:txtbk}

\begin{itemize}
	\item "\textit{TALN - course (ESI)}" [FR].
	Abdelkrime Aries. 
	\kurl{https://github.com/projeduc/ESI_2CS_TALN/releases/download/4.0/taln_v2.pdf}
	
	
	\item "\textit{Speech and Language Processing}".
	Dan Jurafsky and James H. Martin.
	\kurl{https://web.stanford.edu/~jurafsky/slp3/}
	
	\item "\textit{NLP}".
	Jacob Esenstein.
	\kurl{https://cseweb.ucsd.edu/~nnakashole/teaching/eisenstein-nov18.pdf}
	
	\item "\textit{Modern Approaches in NLP - course (LMU Munich)}".
	Carolin Becker, et al.
	\kurl{https://slds-lmu.github.io/seminar_nlp_ss20/index.html}
\end{itemize}

\subsection{To read (ML)}

\begin{itemize}
	\item "\textit{Dive into Deep Learning}" [online].
	Many authors (Amazon).
	\kurl{https://d2l.ai/index.html}
	
	\item "\textit{Deep Learning - course (NYU Center for data science)}".
	Yann LeCun and Alfredo Canziani.
	\kurl{https://atcold.github.io/pytorch-Deep-Learning/}
\end{itemize}

\end{document}
